% Part: first-order-logic
% Chapter: natural-deduction
% Section: proof-theoretic-notions

\documentclass[../../../include/open-logic-section]{subfiles}

\begin{document}

\iftag{FOL}
      {\olfileid{fol}{ntd}{ptn}}
      {\olfileid{pl}{ntd}{ptn}}

\olsection{مفاهيم نظرية البرهان}

\begin{editorial}
في هذا القسم بيان علاقة قابلية الإثبات ومعنى الاتساق في الاستنتاج
الطبيعي.
\end{editorial}

\begin{explain}
قد تقدّم تعريف الصلاحية والاستلزام وقابلية الإرضاء، وهي مفاهيم دلالية.
ولها نظائر من \emph{مفاهيم نظرية البرهان} نبيّنها الآن. وليس مدار هذه
النظائر على إرضاء ال!!{sentence}s في
\iftag{FOL}{ال!!{structure}s}{ال!!{valuation}s}، وإنما مدارها على
!!{derivability} بعض ال!!{sentence}s من غيرها أو !!{nonderivability}ها.
ومن النتائج الجليلة أن المعاني الدلالية ونظائرها البرهانية تتطابق؛
وهذا هو الذي تقرره \emph{مبرهنتا السلامة} و\emph{الاكتمال}.
\end{explain}


\begin{defn}[المبرهنات]
ال!!{sentence}~$!A$ تسمّى \emph{مبرهنةً} متى وجد !!a{derivation} لـ~$!A$
في الاستنتاج الطبيعي، وجميع افتراضاته !!{discharged}. ويرمز إلى ذلك
بـ~$\Proves !A$؛ فإن لم تكن~$!A$ مبرهنة رمزنا إلى ذلك بـ~$\Proves/ !A$.
\end{defn}

\begin{defn}[!!^{derivability}]
نقول إن !!^a{sentence}~$!A$ \emph{!!{derivable} من} مجموعة
!!{sentence}s~$\Gamma$، ونكتب~$\Gamma \Proves !A$، متى وجد
!!{derivation} نتيجته~$!A$، لا يخرج أيّ افتراض فيه عن كونه
!!{discharged} أو من عناصر~$\Gamma$. وأما إذا لم تكن~$!A$
!!{derivable} من~$\Gamma$ فترميز ذلك~$\Gamma \Proves/ !A$.
\end{defn}

\begin{defn}[الاتساق]
مجموعة ال!!{sentence}s~$\Gamma$ \emph{غير متسقة} إذا وفقط إذا
كان~$\Gamma \Proves \lfalse$. فإن انتفى عدم الاتساق عن~$\Gamma$،
أي كان~$\Gamma \Proves/ \lfalse$، سمّيناها \emph{متسقة}.
\end{defn}

\begin{prop}[الانعكاسية]
\ollabel{prop:reflexivity}
إذا كان~$!A \in \Gamma$، فإن~$\Gamma \Proves !A$.
\end{prop}

\begin{proof}
يكفي أن نأخذ الافتراض~$!A$ وحده؛ فهو !!a{derivation} لـ~$!A$، وكل
افتراض !!{undischarged} فيه، وهو~$!A$ نفسه، ينتمي إلى~$\Gamma$.
\end{proof}
  

\begin{prop}[الرتابة]
\ollabel{prop:monotonicity}
إذا كان~$\Gamma \subseteq \Delta$ وكان~$\Gamma \Proves !A$، فإن
$\Delta \Proves !A$.
\end{prop}

\begin{proof}
يكفي !!{derivation} لـ~$!A$ من~$\Gamma$ بعينه؛ إذ هو
!!a{derivation} لـ~$!A$ من~$\Delta$ أيضًا، لدخول المجموعة الأولى في الثانية.
\end{proof}

\begin{prop}[التعدي]
\ollabel{prop:transitivity}
إذا كان~$\Gamma \Proves !A$ وكان~$\{!A\} \cup \Delta \Proves !B$،
فإن~$\Gamma \cup \Delta \Proves !B$.
\end{prop}

\begin{proof}
من~$\Gamma \Proves !A$ نأخذ !!a{derivation}~$\delta_0$ لـ~$!A$،
افتراضاته ال!!{undischarged} كلّها في~$\Gamma$. ومن
$\{!A\} \cup \Delta \Proves !B$ نأخذ !!a{derivation}~$\delta_1$
لـ~$!B$، افتراضاته ال!!{undischarged} كلّها في
$\{!A\} \cup \Delta$. ونركّب الاشتقاقين على الوجه الآتي:
\begin{prooftree}
  \AxiomC{$\Delta, \Discharge{!A}{1}$}
  \RightLabel{$\delta_1$}
  \DeduceC{$!B$}
  \DischargeRule{\Intro{\lif}}{1}
  \UnaryInfC{$!A \lif !B$}
  \AxiomC{$\Gamma$}
  \RightLabel{$\delta_0$}
  \DeduceC{$!A$}
  \RightLabel{\Elim{\lif}}
  \BinaryInfC{$!B$}
\end{prooftree}
فلا يبقى من الافتراضات ال!!{undischarged} إلا ما كان في~$\Gamma \cup
\Delta$، وبذلك ثبت~$\Gamma \cup \Delta \Proves !B$.
\end{proof}

وإذا كانت المجموعة~$\Gamma = \{!A_1, !A_2, \ldots, !A_k\}$ منتهية،
اختصرنا فكتبنا~$!A_1,!A_2,\ldots,!A_k \Proves !B$
مكان~$\Gamma \Proves !B$. ومن ذلك أن قولنا~$!A \Proves !B$
معناه~$\{!A\} \Proves !B$.

ومن التعدي أنه متى كان~$\Gamma \Proves !A$ و$!A \Proves !B$،
لزم~$\Gamma \Proves !B$. وبإجرائه مرارًا، إذا كان
$!A_1, \dots, !A_n \Proves !B$ وكان~$\Gamma \Proves !A_i$ لكل~$i$،
حصل~$\Gamma \Proves !B$.

\begin{prop}
\ollabel{prop:incons}
العبارات الآتية متكافئة:
    \begin{enumerate}
      \item \( \Gamma \) غير متسقة.
      \item \( \Gamma \Proves {!A} \) لكل !!{sentence}~\( {!A} \).
      \item \( \Gamma \Proves {!A} \) و\( \Gamma \Proves \lnot {!A} \) لبعض !!{sentence}~\( {!A} \).
    \end{enumerate}
\end{prop}

\begin{proof}
تمرين.
\end{proof}

\tagprob{FOL}
\begin{prob}
أثبت \olref[fol][ntd][ptn]{prop:incons}.
\end{prob}
\tagendprob

\tagprob{notFOL}
\begin{prob}
أثبت \olref[pl][ntd][ptn]{prop:incons}.
\end{prob}
\tagendprob

\begin{prop}[التراص]
  \ollabel{prop:proves-compact}
  \begin{enumerate}
  \item إذا كان~$\Gamma \Proves !A$، وجدت مجموعة جزئية منتهية
    $\Gamma_0 \subseteq \Gamma$ بحيث~$\Gamma_0 \Proves !A$.
  \item إذا كانت كل مجموعة جزئية منتهية من~$\Gamma$ متسقة، فإن
    $\Gamma$ متسقة.
  \end{enumerate}
\end{prop}

\begin{proof}
  \begin{enumerate}
    \item من~$\Gamma \Proves !A$ نأخذ !!a{derivation}~$\delta$
      لـ~$!A$ من~$\Gamma$، ونسمّي~$\Gamma_0$ مجموعة الافتراضات
      ال!!{undischarged} في~$\delta$. وكل !!{derivation} منته؛ ومن ثم
      لا تضمّ~$\Gamma_0$ إلا عددًا منتهيًا من ال!!{sentence}s.
      فالاشتقاق~$\delta$ نفسه !!a{derivation} لـ~$!A$ من مجموعة
      منتهية~$\Gamma_0 \subseteq \Gamma$.
    \item نأخذ المقابل بالنقيض للفقرة~(1)، مخصّصين إياها بالحالة
      $!A \ident \lfalse$.
  \end{enumerate}
\end{proof}

\end{document}
